\documentclass[12pt,a4paper]{article}

\usepackage[utf8]{inputenc}
\usepackage[T1]{fontenc}
\usepackage[brazil]{babel}
\usepackage[margin=2cm]{geometry}
\usepackage{hyperref}
\usepackage{enumitem}

\title{\textbf{Material de Apoio e Roteiro de Apresentação: Câmbio CVT}}
\author{Lucas Cechin Mário \and André Ghisleni Raimann \\ IFSC — Campus Chapecó}
\date{2026}

\begin{document}
\maketitle

\section*{Instruções para a Apresentação}
Este documento serve como guia de consulta rápida (cola) durante o seminário. Os tópicos aqui descritos expandem os pontos abordados nos slides, fornecendo embasamento mecânico para eventuais perguntas da banca ou da turma.

\section{Introdução e Contexto (Slides 2 e 3)}
\begin{itemize}
    \item \textbf{Como explicar:} Diferente de um câmbio manual ou automático convencional (que tem relações fixas e causa aquelas "trancadas"), o CVT funciona como um rampa lisa. O motor fica cravado na melhor rotação possível enquanto o câmbio adapta a velocidade das rodas.
    \item \textbf{Histórico Rápido:} Mencione Da Vinci como curiosidade (ele pensou no conceito), mas foque que quem colocou para rodar mesmo foi a DAF nos anos 50 com o Variomatic. Hoje, domina os carros japoneses.
\end{itemize}

\section{Fundamentos e Dimensionamento (Slide 4)}
\begin{itemize}
    \item \textbf{Mecânica Pura:} O coração do sistema são as polias cônicas. Quando as faces da polia do motor se aproximam, a correia sobe e aumenta o diâmetro. Enquanto isso, a polia das rodas faz o oposto.
    \item \textbf{Referência Teórica:} Caso alguém pergunte sobre falhas, cite que o dimensionamento à fadiga dessas peças segue os mesmos princípios descritos no \textit{Shigley}. O atrito constante da correia exige eixos e rolamentos superdimensionados.
\end{itemize}

\section{O Diferencial do Seminário: Direct Shift-CVT (Slide 5)}
\begin{itemize}
    \item \textbf{Ponto Forte da Apresentação:} Explique o problema clássico do CVT: ele arrasta para sair do lugar. 
    \item \textbf{A Solução:} A Toyota pegou uma primeira marcha de um câmbio normal e colocou no CVT (Launch Gear). O carro arranca na engrenagem, tira o peso do conjunto de polias (inércia cai 40\%) e só depois joga o movimento pra correia. Melhora a sensação do motorista e salva combustível.
\end{itemize}

\section{Manutenção: A Realidade da Oficina (Slide 6)}
\begin{itemize}
    \item \textbf{Prática de Controle:} Um câmbio CVT vive da precisão hidráulica. Na oficina, não basta "olhar". O diagnóstico (como feito nos JATCO da Nissan/Renault) exige colocar manômetros de mais de 1.000 psi para ver se o LockUp ou as válvulas estão segurando pressão.
    \item \textbf{Sintomas Críticos:} Trepidação ao acelerar quase sempre significa fluido vencido, rolamento fadigado ou ímãs saturados de limalha no cárter (engolindo a malha de controle eletro-hidráulica).
\end{itemize}

\section{Prós e Contras (Slides 7 e 8)}
\begin{itemize}
    \item \textbf{Eficiência e Conforto:} Destaque os bons números de consumo e a vantagem de manter a aceleração constante (bom para ultrapassagens).
    \item \textbf{Drone Noise:} Explique o som de enceradeira/afogamento: como o CVT sobe o giro e trava a rotação lá em cima enquanto a velocidade sobe lentamente, o cérebro humano estranha.
    \item \textbf{Limitações:} Dificilmente você vai ver um CVT em uma caminhonete de 6 toneladas ou em um esportivo puro, o torque esmaga a correia, embora a engenharia moderna já tenha melhorado isso consideravelmente (como o Audi Multitronic).
\end{itemize}

\section{Conclusão}
Fechem frisando a evolução. O CVT não é mais só "aquele câmbio chato de Scooter". Com controle hidráulico moderno e inovações mecânicas (engrenagens acopladas), ele resolveu seus problemas de lentidão e durabilidade e se tornou uma das melhores transmissões do mercado atual em quesito eficiência térmica.

\end{document}